\section{Arithmetic progression discrepancies and atomic prime measures}
\label{acfd:section}

Divisor incidence reduces a prime-weighted quadratic form to a finite
linear functional on progression discrepancies.  The modulus weight in this
functional must be retained when a distribution estimate is applied.
Logarithmic prime measures give a second, distributional realization of
the same counts.  Their topology imposes a separate restriction.

The coefficient spaces and sieve forms are real.  Finite sets carry counting measure.
Integrals on Euclidean spaces use Lebesgue measure, and logarithms are natural.
Write $\mathbb P$ for the positive primes, $\mu$ for the M\"obius function,
$\varphi$ for Euler's totient, and $\omega(n)$ for the number of distinct
prime factors of $n$.  Empty products equal one.

\subsection{The finite arithmetic maps}

Fix distinct nonnegative integers $h_1,\ldots,h_k$, with $k\geq1$.
Assume their set $\mathcal H$ is admissible: the classes $-h_i$ do not
exhaust $\mathbb Z/p\mathbb Z$ for any prime $p$.
Let $W$ be squarefree and contain every prime factor of
$\prod_{i<j}|h_i-h_j|$.  Choose $v$ such that
$\gcd(v+h_i,W)=1$ for all $i$.
Fix integers $N\geq2$ and a real number $R>1$ satisfying $R^2<N$.
Put
\begin{gather*}
 I=\{N\leq n<2N:n\equiv v\pmod W\},\\
 \mathcal D=\{\boldsymbol d\in\mathbb N^k:
     D(\boldsymbol d)<R,\ \mu^2(D(\boldsymbol d))=1,
     \ \gcd(D(\boldsymbol d),W)=1\},
 \qquad D(\boldsymbol d)=\prod_i d_i.
\end{gather*}
The support condition makes the coordinates of each $\boldsymbol d$
pairwise coprime.  Let $Y=\mathbb R^{\mathcal D}$ and
$\Lambda=\mathbb R^{\mathcal D}$.  The maps
$L:Y\to\Lambda$ and $T:\Lambda\to\mathbb R^I$ are
\begin{align}
 (Ly)_{\boldsymbol d}
 &=\left(\prod_i\mu(d_i)d_i\right)
   \sum_{\substack{\boldsymbol r\in\mathcal D\\d_i\mid r_i\ (1\leq i\leq k)}}
       \frac{y_{\boldsymbol r}}{\prod_i\varphi(r_i)},
       \label{acfd:mobius}\\
 (T\lambda)(n)
 &=\sum_{\substack{\boldsymbol d\in\mathcal D\\d_i\mid n+h_i\ (1\leq i\leq k)}}
       \lambda_{\boldsymbol d}.
       \label{acfd:incidence}
\end{align}
For $\lambda,\gamma\in\Lambda$, define
\[
 S_1(\lambda,\gamma)=\sum_{n\in I}(T\lambda)(n)(T\gamma)(n),\qquad
 S_{2,i}(\lambda,\gamma)=\sum_{n\in I}
     \mathbf1_{\mathbb P}(n+h_i)(T\lambda)(n)(T\gamma)(n).
\]
Their sum over $i$ is denoted by $S_2$.  Diagonal values have the
abbreviations $S_1(\lambda)$ and $S_2(\lambda)$.

Set $\Delta_k=\{t\in[0,\infty)^k:\sum_i t_i\leq1\}$ and
$\mathcal X_k=C_c^\infty(\Delta_k^\circ;\mathbb R)$.
Every function is extended by zero to $\mathbb R^k$.
Define
\[
 (PF)_{\boldsymbol r}
   =F\left(\frac{\log r_1}{\log R},\ldots,
           \frac{\log r_k}{\log R}\right),\qquad E=LP,
 \qquad
 \|F\|_{\mathcal X}
   =\sup_{[0,1]^k}\left(|F|+\sum_i|\partial_iF|\right).
\]
Thus $E:\mathcal X_k\to\Lambda$ uses point values on a specified smooth
space.  It does not assign values to arbitrary $L^2$ equivalence classes.

\begin{lemma}[Support and coefficient bounds]
Let\label{acfd:coefficients} $F\in\mathcal X_k$.  Suppose its support lies in
$\{\sum_i t_i\leq1-\sigma\}$, where $0<\sigma<1$.
Then $EF$ vanishes unless $D(\boldsymbol d)\leq R^{1-\sigma}$.
With $\ell=1+\log R$, one has
\[
 \|EF\|_{\ell^\infty(\mathcal D)}
       \leq e^{k+1}\ell^{k+1}\|F\|_\infty,
 \qquad
 \|EF\|_{\ell^1(\mathcal D)}
       \leq e^{k+1}R^{1-\sigma}\ell^{2k}\|F\|_\infty.
\]
The same estimates with $\sigma=0$ hold without a specified support margin.
\end{lemma}

\begin{proof}
A nonzero sample has $D(\boldsymbol r)\leq R^{1-\sigma}$.
Every contributing $\boldsymbol d$ divides such a $\boldsymbol r$
coordinatewise, which proves the support assertion.
For $X\geq1$, the identity
\[
 \frac n{\varphi(n)}=\sum_{d\mid n}\frac{\mu^2(d)}{\varphi(d)}
\]
follows by checking prime powers and multiplying.  Therefore
\begin{equation}\label{acfd:totient-sum}
 \sum_{n\leq X}\frac1{\varphi(n)}
 \leq(1+\log X)\sum_{d\geq1}\frac{\mu^2(d)}{d\varphi(d)}
 \leq e(1+\log X).
\end{equation}
Indeed, finite Euler products for the last sum are bounded by
$\exp(\sum_{p}1/[p(p-1)])\leq e$.
Monotone convergence proves the bound for the infinite sum.

Write $r_i=d_is_i$ in \eqref{acfd:mobius}.
Squarefreeness gives $\varphi(r_i)=\varphi(d_i)\varphi(s_i)$.
Dropping the remaining restrictions on the $s_i$ bounds their sum by
$(e\ell)^k$.
Also $\prod_i d_i/\varphi(d_i)=D(\boldsymbol d)/\varphi(D(\boldsymbol d))
\leq e\ell$, by the preceding identity and estimate.
This proves the supremum bound.
For $X\geq1$, the number of ordered positive integer tuples with product
at most $X$ is at most $X(1+\log X)^{k-1}$.
Fixing the first $k-1$ coordinates proves this by summing
$X/(d_1\cdots d_{k-1})$ and enlarging each remaining range to $[1,X]$.
Apply this bound with $X=R^{1-\sigma}$ to obtain the $\ell^1$ estimate.
\end{proof}

\subsection{Exact progression discrepancies}

For $\boldsymbol d,\boldsymbol e\in\mathcal D$, write
$q_j=[d_j,e_j]$ for their least common multiple.
Call the pair compatible when the $q_j$ are pairwise coprime.
An incompatible pair contributes zero to every form above.
In fact, a common prime divisor of $q_j$ and $q_l$ would divide
$h_j-h_l$ and hence $W$, contrary to the definition of $\mathcal D$.

For a compatible pair with $d_i=e_i=1$, put
\[
 Q_i(\boldsymbol d,\boldsymbol e)=W\prod_{j\ne i}q_j.
\]
There is a unique reduced residue $a_i$ modulo $Q_i$ satisfying
\[
 a_i\equiv v+h_i\pmod W,\qquad
 a_i\equiv h_i-h_j\pmod {q_j}\quad(j\ne i).
\]
These residues are reduced because all prime divisors of $h_i-h_j$
belong to $W$.  The Chinese remainder theorem supplies existence and uniqueness.
When $Q_i=1$, its single residue is treated as reduced and $\varphi(1)=1$.

Let
\begin{align*}
 \Pi_i&=\#\{p\in\mathbb P:N+h_i\leq p<2N+h_i\},\\
 \mathcal E_i(Q,a)
   &=\#\{p\in\mathbb P:N+h_i\leq p<2N+h_i,
                               \ p\equiv a\pmod Q\}
                   -\frac{\Pi_i}{\varphi(Q)},\\
 \Delta_i(Q)&=\max_{\gcd(a,Q)=1}|\mathcal E_i(Q,a)|.
\end{align*}
These are exact interval discrepancies.  No prime asymptotic enters their definition.

\begin{theorem}[The finite discrepancy functional]\label{acfd:exact}
For $\lambda,\gamma\in\Lambda$, define
\[
 B_i^{\mathrm{loc}}(\lambda,\gamma)
   =\sum_{\substack{\boldsymbol d,\boldsymbol e\text{ compatible}\\d_i=e_i=1}}
        \frac{\lambda_{\boldsymbol d}\gamma_{\boldsymbol e}}{\varphi(Q_i)},
 \qquad
 A^{\mathrm{loc}}(\lambda,\gamma)
   =\frac NW\sum_{\boldsymbol d,\boldsymbol e\text{ compatible}}
        \frac{\lambda_{\boldsymbol d}\gamma_{\boldsymbol e}}{\prod_jq_j}.
\]
Then
\begin{align}
 S_{2,i}(\lambda,\gamma)
   &=\Pi_i B_i^{\mathrm{loc}}(\lambda,\gamma)
     +\sum_{\substack{\boldsymbol d,\boldsymbol e\text{ compatible}\\d_i=e_i=1}}
         \lambda_{\boldsymbol d}\gamma_{\boldsymbol e}
         \mathcal E_i(Q_i,a_i),\label{acfd:prime-spine}\\
 |S_1(\lambda,\gamma)-A^{\mathrm{loc}}(\lambda,\gamma)|
   &\leq\|\lambda\|_1\|\gamma\|_1.\label{acfd:mass-spine}
\end{align}
Suppose both coefficient vectors vanish for $D(\boldsymbol d)>X$, with
$1\leq X\leq R$.  Set $z=\max\{1,3(k-1)\}$.
The remainder in \eqref{acfd:prime-spine} satisfies
\begin{equation}\label{acfd:weighted-error}
 |\mathcal R_i(\lambda,\gamma)|
 \leq\|\lambda\|_\infty\|\gamma\|_\infty
      \sum_{Q\leq WX^2}\mu^2(Q)z^{\omega(Q)}\Delta_i(Q).
\end{equation}
\end{theorem}

\begin{proof}
Expand the incidence products and interchange finite sums.
For a compatible pair, the unweighted inner sum counts one residue class
modulo $W\prod_jq_j$ in $N$ consecutive integers.
Its distance from $N/(W\prod_jq_j)$ is at most one.
Taking absolute values proves \eqref{acfd:mass-spine}.

In the $i$th prime sum, $q_i$ divides the prime $n+h_i\geq N$.
But $q_i\leq D(\boldsymbol d)D(\boldsymbol e)<R^2<N$.
Consequently $q_i=1$, or equivalently $d_i=e_i=1$.
The remaining conditions are precisely $n+h_i\equiv a_i\pmod {Q_i}$.
Subtracting its defining mean gives \eqref{acfd:prime-spine}.

For fixed $Q$, each prime dividing $Q/W$ chooses a coordinate $j\ne i$.
It then occurs in $d_j$ only, in $e_j$ only, or in both.
Thus at most $[3(k-1)]^{\omega(Q/W)}$ ordered pairs give $Q$, when $k\geq2$.
The product restrictions can only reduce this count.
Every modulus is squarefree and at most $WX^2$.
Since $z^{\omega(Q/W)}\leq z^{\omega(Q)}$, taking absolute values gives the bound.
For $k=1$, the sole surviving pair is $d_1=e_1=1$ and $Q=W$.
The choice $z=1$ includes this case.
\end{proof}

Equivalently, let $\mathfrak A_i$ be the finite set of reduced pairs $(Q,a)$
which occur above.  Define a bilinear map
\[
 \mathcal C_i:\Lambda\times\Lambda\longrightarrow\mathbb R^{\mathfrak A_i},
 \qquad
 \mathcal C_i(\lambda,\gamma)_{Q,a}
    =\sum_{(Q_i,a_i)=(Q,a)}\lambda_{\boldsymbol d}\gamma_{\boldsymbol e}.
\]
The remainder is exactly the dual pairing
$\langle\mathcal C_i(\lambda,\gamma),\mathcal E_i\rangle$.
This pairing specifies the arithmetic object needed after the finite Rayleigh principle.

\begin{example}[One residue class]
Take\label{acfd:first-model} $k=1$, $h_1=0$, $N=30$, $W=2$, $v=1$, and $R=5$.
Then $\mathcal D=\{1,3\}$.  Let $y=(1,1)$, so
$\lambda=Ly=(3/2,-3/2)$.
Among the fifteen odd integers in $[30,60)$, five are divisible by three.
Therefore $S_1(\lambda)=10(3/2)^2=45/2$.
The seven primes are $31,37,41,43,47,53,59$.
Only $d=e=1$ survives in the prime sum, giving
$S_2(\lambda)=7(3/2)^2=63/4$.
Here $\mathcal E_1(2,1)=0$ and $S_2/S_1=7/10$.
The unrestricted coefficient vector $y$ in this example need not be a
sample of an interior smooth function.
\end{example}

\subsection{From an unweighted distribution hypothesis to the weighted error}

Fix $0<\vartheta\leq1$ and assume that, for every $A>0$,
\begin{equation}\label{acfd:lod}
 \sum_{q\leq x^\vartheta}\max_{\gcd(a,q)=1}
 \left|\pi(x;q,a)-\frac{\pi(x)}{\varphi(q)}\right|
       \leq C_A\frac{x}{(\log x)^A}
       \qquad(x\geq x_A).
\end{equation}
The constants may depend on $A$ and $\vartheta$.
This is the distribution convention of \eqref{eq:lod}, as in
\cite[(1.3), p.~384]{Maynard}.
The hypothesis remains an assumption at the chosen exponent.

Fix $0<\delta<\vartheta/2$ and put
\[
 \rho=\vartheta/2-\delta,\qquad
 R_N=N^\rho,\qquad D_0=\log\log\log N,\qquad
 W_N=\prod_{p\leq D_0}p.
\]
Use these parameters in the finite maps for sufficiently large integer $N$.
The shifts remain fixed.  Thus $W_N$ eventually contains the prime factors
of every shift difference, and an admissible class $v_N$ exists.

\begin{theorem}[Uniform control of the arithmetic remainder]\label{acfd:lod-transfer}
Assume \eqref{acfd:lod}.  For every $B>0$, there are constants
$C_B'$ and $N_B$ independent of $F,G\in\mathcal X_k$ and of $v_N$ such that
\begin{equation}\label{acfd:remainder-decay}
 |\mathcal R_i(E_NF,E_NG)|
       \leq C_B' N(\log N)^{-B}\|F\|_\infty\|G\|_\infty
       \qquad(N\geq N_B).
\end{equation}
If both smooth supports have slack at least $\sigma>0$, every modulus
in the exact discrepancy functional is at most
$W_NR_N^{2(1-\sigma)}$.
Without that margin, all moduli are less than $W_NR_N^2$.
\end{theorem}

\begin{proof}
The elementary bound $W_N\leq D_0^{D_0}$ gives $W_N=N^{o(1)}$.
Consequently
\[
 W_NR_N^2\leq N^{\vartheta-\delta}<N
\]
for all sufficiently large $N$.
For the two integer endpoints $x_+=2N+h_i-1$ and $x_-=N+h_i-1$,
the same bound is eventually at most $x_\pm^\vartheta$.
The interval discrepancy is the difference of the two endpoint discrepancies.
Applying \eqref{acfd:lod} at both endpoints yields
\[
 \sum_{Q\leq W_NR_N^2}\Delta_i(Q)
        \ll_A N(\log N)^{-A}.
\]
This argument preserves the half-open interval and its fixed shift.

For $Q\leq N$, the trivial counting bound gives
\[
 \Delta_i(Q)\leq N/Q+1+N/\varphi(Q)\leq3N/\varphi(Q).
\]
Let $b=z^2$, an integer.  For any $Y\geq1$,
\begin{equation}\label{acfd:moment}
 \sum_{Q\leq Y}\frac{\mu^2(Q)b^{\omega(Q)}}{\varphi(Q)}
       \leq\left(\sum_{n\leq Y}\frac{\mu^2(n)}{\varphi(n)}\right)^b
       \leq e^b(1+\log Y)^b.
\end{equation}
To see the first inequality, expand the $b$th power.
Each squarefree $Q$ has $b^{\omega(Q)}$ allocations of its primes to
$b$ ordered factors.  All these factors are at most $Y$.
Their denominators multiply to $\varphi(Q)$, and other summands are nonnegative.

Cauchy--Schwarz, the trivial discrepancy bound, and
\eqref{acfd:moment} now give
\begin{align*}
 \sum_{Q\leq W_NR_N^2}\mu^2(Q)z^{\omega(Q)}\Delta_i(Q)
 &\leq
 \left(\sum_{Q\leq W_NR_N^2}\Delta_i(Q)\right)^{1/2}
 \left(3N\sum_{Q\leq W_NR_N^2}
               \frac{\mu^2(Q)z^{2\omega(Q)}}{\varphi(Q)}\right)^{1/2}\\
 &\ll_A N(\log N)^{(b-A)/2}.
\end{align*}
Lemma~\ref{acfd:coefficients} contributes a further factor
$O_k((\log N)^{2k+2})\|F\|_\infty\|G\|_\infty$.
Choose $A\geq b+4k+4+2B$ and apply \eqref{acfd:weighted-error}.
This proves \eqref{acfd:remainder-decay}.
The support assertion follows from Lemma~\ref{acfd:coefficients}.
\end{proof}

The proof displays the weighted distribution step used in Maynard's
equations (5.18)--(5.20), p.~396.  A bound for an unweighted sum of
discrepancies cannot simply discard the multiplicity $z^{\omega(Q)}$.
The positive second moment in \eqref{acfd:moment} is the required additional estimate.

\subsection{The remaining comparison with the simplex forms}

Let $H_k=L^2(\Delta_k,dt;\mathbb R)$.
For $1\leq i\leq k$, define
\[
 (T_iF)(t_{\widehat i})=
    \int_0^{1-\sum_{j\ne i}t_j}F(t_1,\ldots,t_k)\,dt_i,
 \qquad K_k=\sum_iT_i^*T_i.
\]
Here $T_i:H_k\to L^2(\Delta_{k-1})$ uses ordinary length on each fiber.
For $k=1$, $\Delta_0$ is one point of mass one.
Fiber Cauchy--Schwarz shows $\|T_i\|\leq1$ and $0\leq K_k\leq kI$.
Set
\[
 C_N=\frac{\varphi(W_N)^kN(\log R_N)^k}{W_N^{k+1}}.
\]
The arithmetic remainder theorem compares exact prime sums to finite
totient sums.  The following explicit hypotheses state the further comparison
needed to reach $K_k$:
\begin{align}
 |A^{\mathrm{loc}}(E_NF,E_NG)-C_N\langle F,G\rangle|
     &\leq C_N\alpha_N\|F\|_{\mathcal X}\|G\|_{\mathcal X},
       \label{acfd:mass-local-hyp}\\
 \left|\sum_i\Pi_iB_i^{\mathrm{loc}}(E_NF,E_NG)
                -\rho C_N\langle F,K_kG\rangle\right|
     &\leq C_N\beta_N\|F\|_{\mathcal X}\|G\|_{\mathcal X},
       \label{acfd:prime-local-hyp}
\end{align}
where $\alpha_N,\beta_N\geq0$ tend to zero.
These are hypotheses about the displayed finite sums, with uniform quantifiers.

\begin{proposition}[Composition of the comparisons]
Assume\label{acfd:composition} \eqref{acfd:lod}, \eqref{acfd:mass-local-hyp}, and
\eqref{acfd:prime-local-hyp}.  There are nonnegative sequences
$u_N,v_N\to0$ such that, for all $F,G\in\mathcal X_k$,
\begin{align*}
 |S_1(E_NF,E_NG)/C_N-\langle F,G\rangle|
     &\leq(\alpha_N+u_N)\|F\|_{\mathcal X}\|G\|_{\mathcal X},\\
 |S_2(E_NF,E_NG)/(\rho C_N)-\langle F,K_kG\rangle|
     &\leq\rho^{-1}(\beta_N+v_N)
                   \|F\|_{\mathcal X}\|G\|_{\mathcal X}.
\end{align*}
One may take
\[
 u_N=\frac{e^{2k+2}R_N^2(1+\log R_N)^{4k}}{C_N},\qquad
 v_N=\frac{kC_B'N(\log N)^{-B}}{C_N}
\]
for any fixed $B>1$ and all sufficiently large $N$.
In particular, if a fixed nonzero $F$ satisfies
$\rho\langle F,K_kF\rangle>m\|F\|_2^2$, with $m\in\mathbb Z_{\geq0}$,
then some $n\in I_N$ has at least $m+1$ prime translates for all large $N$.
\end{proposition}

\begin{proof}
Theorem~\ref{acfd:exact}, Lemma~\ref{acfd:coefficients}, and
Theorem~\ref{acfd:lod-transfer} give the two bounds by addition.
To verify the limits, note that
\[
 C_N\geq
 \frac{N(\rho\log N)^k}{e^kW_N(1+\log W_N)^k}.
\]
Moreover, $W_N(1+\log W_N)^k\leq\log N$ eventually.
This follows from $\log W_N\leq D_0\log D_0$
and $D_0=\log\log\log N$.
Thus $C_N\geq e^{-k}\rho^kN(\log N)^{k-1}$ eventually.
Since $2\rho<1$, the displayed $u_N,v_N$ tend to zero.

For the fixed function $F$, the two form limits give
$S_1(E_NF)>0$ and $S_2(E_NF)-mS_1(E_NF)>0$ eventually.
If every prime count were at most $m$, each nonnegative squared incidence
weight would give the reverse inequality.  Integrality therefore supplies
at least $m+1$ prime translates.
\end{proof}

For the smooth sieve weights, Maynard's Lemmas 6.2--6.3, pp.~401--404,
give additive comparisons for the exact forms in the norm $\|\cdot\|_{\mathcal X}$.
Those estimates also evaluate the local main sums above, after subtracting
the remainders already bounded here.  Their use retains the hypothesis
\eqref{acfd:lod}.  The finite discrepancy identity alone does not evaluate
either local main sum, supply a distribution exponent, or construct a map
on the full coefficient quotient.

\subsection{Logarithmic prime distributions}

Let $\mathcal D(\mathbb R)=C_c^\infty(\mathbb R;\mathbb R)$ have its usual
test-function topology, and let $\mathcal D'(\mathbb R)$ be its continuous dual.
For $\sigma\in\mathbb R$, define positive atomic measures
\begin{equation}\label{acfd:prime-measures}
 \nu_\sigma=\sum_{p\in\mathbb P}(\log p)p^{-\sigma}\delta_{\log p},
 \qquad
 \mu_\sigma=\sum_{p\in\mathbb P}\sum_{m\geq1}
                  (\log p)p^{-m\sigma}\delta_{m\log p}.
\end{equation}
The second measure is the logarithmic image of the prime-power measure
$\sum_{n\geq2}\Lambda(n)n^{-\sigma}\delta_n$.
Here $\Lambda(p^m)=\log p$ and $\Lambda(n)=0$ otherwise.
This notation $\Lambda(n)$ denotes the von Mangoldt function, distinct from
the finite vector space $\Lambda$ used above.
Every compact interval meets only finitely many atoms.
Thus both formulas define distributions of order zero for every real $\sigma$.

\begin{theorem}[The exact temperateness boundary]
The\label{acfd:tempered} positive measures $\nu_\sigma$ and $\mu_\sigma$ extend continuously to
the Schwartz space $\mathcal S(\mathbb R)$ if and only if $\sigma\geq1$.
For $\sigma\geq1$,
\[
 |\nu_\sigma(f)|\leq\mu_\sigma(|f|)
 \leq e\left(\sum_{j\geq1}j^{-2}\right)
               \sup_{t\in\mathbb R}(1+|t|)^3|f(t)|.
\]
For arbitrary $\sigma\in\mathbb R$ and $T\geq0$,
\begin{equation}\label{acfd:compact-mass}
 \mu_\sigma([0,T])
       \leq e\,e^{\max(1-\sigma,0)T}(T+2)^2.
\end{equation}
\end{theorem}

\begin{proof}
For integers $j\geq0$, there are at most $e^{j+1}$ positive integers
in $[e^j,e^{j+1})$.  At any prime power in this interval,
$\Lambda(n)/n\leq(j+1)e^{-j}$.
Consequently $\mu_1([j,j+1))\leq e(j+1)$.
Summing this bound against $(1+j)^{-3}$ proves the displayed Schwartz estimate.
For $\sigma\geq1$, each mass decreases, so the same bound applies.
Summing without a test function gives
$\mu_1([0,T])\leq e(T+2)^2$.
The identity $d\mu_\sigma(t)=e^{(1-\sigma)t}d\mu_1(t)$ on its support
then proves \eqref{acfd:compact-mass}.

For the converse, first recall an elementary divergence.
If $\sum_p1/p$ were finite, the products
$\prod_{p\leq X}(1-1/p)^{-1}$ would be bounded, since
$-\log(1-1/p)\leq2/p$.
Their geometric-series expansions contain every term $1/n$ with $n\leq X$.
This contradicts divergence of the harmonic series.  Hence $\sum_p1/p=\infty$.

Suppose $\sigma<1$ and choose $a=(1-\sigma)/2>0$.
There is a nonnegative smooth function $f$ which vanishes for $t\leq0$
and equals $e^{-at}$ for $t\geq\log2$.
This function lies in $\mathcal S(\mathbb R)$.
Its putative positive pairing satisfies
\[
 \nu_\sigma(f)
   =\sum_p(\log p)p^{-(1+\sigma)/2}=+\infty,
\]
because each summand is eventually at least $1/p$.
For rigor, choose $0\leq\eta\leq1$ in $C_c^\infty(\mathbb R)$,
equal to one on $[-1,1]$, and set $f_L(t)=f(t)\eta(t/L)$.
Then $f_L\to f$ in every Schwartz seminorm, by exponential decay and
the product rule.  Yet $\nu_\sigma(f_L)\to+\infty$:
every finite positive partial sum is included when $L$ is sufficiently large.
Continuity of a tempered extension would force finite convergence.
This contradiction also applies to $\mu_\sigma$, which dominates $\nu_\sigma$.
\end{proof}

\subsection{Recovering the progression counts by smooth tests}

The finite counts have an exact realization on the distributional carrier.
It also exhibits the derivative cost of separating adjacent integers after logarithmic change of variables.
Fix $\sigma\geq0$ and $H=\max_i h_i$.
Choose $\chi\in C_c^\infty((-1,1))$ with $\chi(0)=1$, and put
$r_N=[4(2N+H+1)]^{-1}$.
For every integer $n\in[N,2N+H)$, define
\[
 b_{n,N}(t)=\frac{n^\sigma}{\log n}\,
                 \chi\left(\frac{t-\log n}{r_N}\right).
\]
For a residue $a$ modulo $Q$, set
\begin{align*}
 \Phi_{i,Q,a,N}(t)
   &=\sum_{N+h_i\leq n<2N+h_i}
          \left(\mathbf1_{n\equiv a\ (Q)}-\frac1{\varphi(Q)}\right)b_{n,N}(t),\\
 \Psi_{i,Q,a,N}(t)
   &=\sum_{\substack{N+h_i\leq n<2N+h_i\\n\equiv a\ (Q)}}b_{n,N}(t).
\end{align*}

\begin{proposition}[Exact logarithmic realization]
For\label{acfd:log-realization} every reduced residue $a$ modulo $Q$,
\[
 \nu_\sigma(\Phi_{i,Q,a,N})=\mathcal E_i(Q,a),\qquad
 \nu_\sigma(\Psi_{i,Q,a,N})
    =\#\{p\in[N+h_i,2N+h_i):p\equiv a\pmod Q\}.
\]
Both tests are supported in
$[\log N-r_N,\log(2N+H)+r_N]$.
For each derivative order $j\geq0$, either test $\Theta$ satisfies
\begin{equation}\label{acfd:test-norm}
 \|\Theta^{(j)}\|_\infty
    \leq\frac{(2N+H)^\sigma}{\log N}
          [4(2N+H+1)]^j\|\chi^{(j)}\|_\infty.
\end{equation}
\end{proposition}

\begin{proof}
For consecutive integers in the indicated range, including its adjacent endpoints,
\[
 \log(n+1)-\log n\geq\frac1{n+1}\geq\frac1{2N+H+1}=4r_N.
\]
The bump supports are therefore disjoint and contain no other logarithm of an integer.
At $t=\log n$, the product of the prime mass and the bump value is one
if $n$ is prime.  It is zero if $n$ is composite, because $\nu_\sigma$
has no atom there.  Finite summation proves both identities.
The coefficients in parentheses have absolute value at most one.
Disjoint supports and differentiation of a rescaled bump prove
\eqref{acfd:test-norm}.
\end{proof}

Combining Proposition~\ref{acfd:log-realization} with
Theorem~\ref{acfd:exact} expresses the arithmetic remainder as
\[
 \mathcal R_i(\lambda,\gamma)
   =\nu_\sigma\!\left(
      \sum_{\substack{\boldsymbol d,\boldsymbol e\text{ compatible}\\d_i=e_i=1}}
         \lambda_{\boldsymbol d}\gamma_{\boldsymbol e}
         \Phi_{i,Q_i,a_i,N}\right).
\]
This is an equality of finite sums on compact smooth tests.
Its derivative bound grows with $N$; distributional continuity alone supplies
no decay for the arithmetic remainder.

At $\sigma=1/2$, the unrenormalized prime-power term
\[
 f\longmapsto\sum_{p,m\geq1}(\log p)p^{-m/2}
                       \bigl(f(m\log p)+f(-m\log p)\bigr)
\]
also fails to extend to the full Schwartz space.
Apply the divergence proof to a nonnegative even test.
It is well defined on compact smooth tests, with the bound
\eqref{acfd:compact-mass} for both signs.
A comparison with a global explicit formula must specify its test space,
the remaining terms, their normalization, and any cancellation or subtraction.
The obstruction concerns this positive term by itself; it does not exclude
a separately constructed distribution for a combined formula.

Finally, the sampling map has its own distinct limitation.
For fixed $N$, a nonzero smooth bump can avoid every point of the finite
set $x_N(\mathcal D_N)$.
Normalizing that bump in $H_k$ gives $\|F_N\|_2=1$ and $E_NF_N=0$.
At any interior sample, smooth bumps of fixed value and shrinking support
have $L^2$ norms tending to zero.
Thus point sampling cannot extend to a bounded map on $H_k$ when such a node exists.
Neither the finite logarithmic realization nor the smooth form comparison
constructs a uniform reconstruction of all coefficient classes modulo $\ker T_N$.
That reconstruction requires its own norm estimates and comparison with $K_k$.
